% ============================================================================

% Paper: Frequent Losers as Cover Bidders in Public Procurement
% ============================================================================

\section{Empirical Strategy}
\label{sec:empirical}

The empirical analysis has three parts that play distinct roles.
First, we estimate the conditional price association between FL
presence and tender outcomes using OLS with high-dimensional fixed
effects, temporal cross-fitting, and matching
(Section~\ref{sec:reduced_form}); these document a robust
association, not a causal effect.
Second, we evaluate the screen's detection performance against CADE
cartel convictions and benchmark it against bid-level screens
(Section~\ref{sec:detection_validation}); this is the paper's
primary validation exercise.
Third, we report diagnostics assessing whether the screen's behavior
is coherent with coordinated cover bidding, drawing on structural
estimation (Section~\ref{sec:structural_estimation}), Bajari--Ye
bid-coordination tests (Section~\ref{sec:bajari_ye}), and
network-split heterogeneity
(Section~\ref{sec:empirical_network}); these support the screen's
economic rationale without claiming mechanism identification.

% ── 5.1 Reduced-Form Identification ──────────────────────────────

\subsection{Conditional Price Association}
\label{sec:reduced_form}

\paragraph{OLS baseline.}
The baseline specification regresses tender outcomes on FL presence:
\begin{equation}
  y_{igt} = \beta \cdot \text{losers}_{igt} + \mathbf{x}_{igt}'
  \boldsymbol{\delta} + \alpha_g + \lambda_t + \gamma_k + \varepsilon_{igt}
  \label{eq:ols_main}
\end{equation}
where $y_{igt}$ is the outcome for tender-item $i$ in item group $g$ at
time $t$ and purchasing unit $k$; $\text{losers}_{igt} \in \{0,1\}$
indicates FL presence; $\mathbf{x}_{igt}$ includes a convite indicator;
$\alpha_g$, $\lambda_t$, and $\gamma_k$ are item, year, and PBU fixed
effects; and $\varepsilon_{igt}$ is the error term clustered at the item
level.

The OLS coefficient $\hat{\beta}$ captures a conditional
association---the average difference in outcomes between FL-present and
FL-absent tenders within the same item type, year, and purchasing unit.
Under the model's first prediction (Table~\ref{tab:predictions}), $\hat{\beta} > 0$ for prices.
Item, year, and PBU fixed effects absorb time-invariant cross-market
heterogeneity; the residual threat is time-varying within-market
confounders (Section~\ref{sec:limitations}). We assess the sensitivity
of $\hat{\beta}$ to such confounders using the
\citet{cinelli2020making} framework, which yields a robustness value
of $RV_{q=1} = 17.5\%$: an unobserved confounder would need to
explain at least 17.5\% of the residual variation in both FL presence
and log prices to reduce $\hat{\beta}$ to zero.

\paragraph{Matching estimators.}
To address potential selection on observables, we estimate the
price association using two matching approaches: coarsened exact
matching (CEM) on year, procedure type, item group, and PBU size
quartile ($\hat{\beta}_{\text{CEM}} = 0.077$, $N = 969{,}751$), and
inverse probability weighting (IPW, $\hat{\beta}_{\text{IPW}} = 0.055$,
$N = 830{,}194$). These estimates bracket the OLS baseline (0.064)
and the cross-fit average (0.036), providing a non-parametric check
that the price association survives reweighting on observables.

\paragraph{Instrumental variable (measurement-error diagnostic).}
The binary FL indicator contains misclassification noise: firms near
the IQR threshold may not be cover bidders, and some cover bidders
escape the threshold. Classical binary measurement error attenuates
the OLS coefficient toward zero. We use a leave-one-out FL supply
instrument primarily to diagnose this attenuation:
\begin{equation}
  Z_{kgt} = \sum_{j \neq k} \mathbf{1}[\text{FL firm active at PBU } j
  \text{ in group } g, \text{ year } t]
  \label{eq:iv_v6}
\end{equation}
The instrument exploits aggregate FL supply variation at other
purchasing units in the same product market and year. We treat
$\hat{\beta}_{\text{IV}}$ as a measurement-error diagnostic: it
indicates the direction of bias in the binary FL indicator
(attenuation toward zero) but not the true magnitude, because
balance tests reveal systematic imbalance across observables
(Table~\ref{tab:iv_balance}, Appendix~\ref{app:robustness}). The IV
is consistent with OLS understating the association; our primary range is
cross-fit (0.036), OLS (0.064), and matching (0.055--0.077).

% ── 5.2 Detection Validation ─────────────────────────────────────

\subsection{Detection Performance Validation}
\label{sec:detection_validation}

The detection assessment uses CADE cartel convictions as ground truth,
with FL status defined on 2009--2014 data and evaluated against
2015--2019 enforcement outcomes (the temporal split described in
Section~\ref{sec:cade_sample}).

\paragraph{ROC analysis.}
We vary the IQR multiplier from 0.5 to 5.0 in steps of 0.1, producing
46 FL thresholds. For each, we compute the true-positive and
false-positive rates against CADE co-participation and trace the ROC
curve. The optimal multiplier maximizes Youden's $J$ (TPR $-$ FPR),
connecting the IQR rule to a formal detection-theoretic criterion.

\paragraph{Benchmarking.}
We benchmark FL against \citet{imhof2019detecting}-style bid-level
features on the same CADE ground truth, and include both screens in a
horse-race regression to test orthogonality
(Section~\ref{sec:results_detection}).

% ── 5.3 Structural Estimation ────────────────────────────────────

\subsection{Structural Likelihood Estimation}
\label{sec:structural_estimation}

The mixture model (Appendix~\ref{app:structural_id}) is estimated in two stages on bid-level data. After restricting to losing bids with valid prices and item-year fixed-effect coverage, the estimation sample comprises 23.2 million genuine (non-FL) losing bids and 189,381 FL losing bids (Table~\ref{tab:structural_params}).

\paragraph{Stage 1: Genuine bid distribution.}
We estimate $(\boldsymbol{\beta}, \sigma_g^2)$ from the 23.2 million non-FL losing bids by maximum likelihood under a log-normal specification:
\begin{equation}
  (\hat{\boldsymbol{\beta}}, \hat{\sigma}_g^2) =
  \operatorname*{arg\,max}_{\boldsymbol{\beta}, \sigma_g^2}
  \sum_{t=1}^{T} \sum_{j \in \mathcal{G}_t}
  \log \phi\!\left(
    \frac{\log b_{jt} - \mathbf{x}_j'\boldsymbol{\beta}
    - \hat{\alpha}_i - \hat{\lambda}_t}{\sigma_g}
  \right)
  \label{eq:mle_genuine}
\end{equation}
where $\phi(\cdot)$ is the standard normal density and $\hat{\alpha}_i$, $\hat{\lambda}_t$ are absorbed item and year fixed effects. In practice, this reduces to OLS of $\log b_{jt}$ on $\mathbf{x}_j$ with fixed effects, yielding $\hat{\sigma}_g^2$ as the residual variance.

\paragraph{Stage 2: Cover bid distribution.}
Conditional on the winning bid $b^*_t$, we estimate cover-bid parameters from FL bids under each regime specification (Appendix~\ref{app:structural_id}).

\paragraph{Model selection.}
We select between regimes using BIC; the implied markup is:
\begin{equation}
  \hat{\mu}_{\text{structural}} =
  \frac{E[b^*_t \mid \text{FL present}] -
        E[b^*_t \mid \text{FL absent}]}{E[b^*_t \mid \text{FL absent}]}
  \label{eq:structural_markup}
\end{equation}
We report this alongside the reduced-form estimates. Standard errors
are computed via parametric bootstrap (500 replications, tender-level
resampling).

% ── 5.4 Bajari-Ye Tests ──────────────────────────────────────────

\subsection{Bajari--Ye Bid Coordination Tests}
\label{sec:bajari_ye}

We implement the \citet{bajari2003deciding} tests using corrected
first-stage residuals from the bid equation:
\begin{equation}
  \log b_{jt} = \mathbf{x}_j'\boldsymbol{\phi} + \alpha_i + \lambda_t
  + u_{jt}
  \label{eq:by_first_stage}
\end{equation}
where $\mathbf{x}_j$ includes firm size, firm age, and CNAE sector
dummies. We exclude $n_{\text{bids}}$ from the first stage: it is
jointly determined with bid levels and mechanically increased by FL
presence \citep[p.~978]{bajari2003deciding}.

Using the residuals $\hat{u}_{jt}$, we test (i)~\emph{exchangeability}
(KS test comparing FL and non-FL residual distributions;
Diagnostic~\ref{pred:exchange_v6}) and (ii)~\emph{conditional
independence} (mean pairwise product of FL residuals within tenders;
bootstrap with 1,000 iterations). A fake-groups placebo splits non-FL
firms randomly within each tender; the economically relevant statistic
is the FL--non-FL pairwise \emph{difference}, which isolates excess
FL correlation beyond common tender-level shocks. We also report a
tender-FE variant as a conservative specification absorbing all
tender-level variation.

% ── 5.5 Network-Split Heterogeneity ──────────────────────────────

\subsection{Network-Split Heterogeneity Test}
\label{sec:empirical_network}

To test Proposition~\ref{prop:market_selection}---that cover bidding
concentrates in competitive markets---we split FL firms into two
subgroups based on co-bidding network characteristics.

\paragraph{Winner HHI classification.}
For each FL firm $j$, we compute the HHI of winning firms across all
tenders in which $j$ participates. Firms above the sample-median HHI
are classified as \emph{concentrated-market}; those below as
\emph{competitive-market}. We interact $\text{losers}_{igt}$ with the
market-structure indicator in Equation~\eqref{eq:ols_main}.

Before presenting the price regressions, we check the FL
classification against external enforcement data
(Section~\ref{sec:cade}). This provides an initial check that FL presence flags
economically meaningful environments---a necessary condition for the
price association to be informative.
